\documentclass[10pt]{extarticle}
\usepackage[a4paper,margin=1.5cm]{geometry}
\usepackage[T1]{fontenc}
\usepackage{amsmath,amssymb}
\usepackage{graphicx}
\usepackage[export]{adjustbox}
\usepackage{multicol}
\usepackage{enumitem}
\usepackage{booktabs,array,tabularx}
\setlength{\columnsep}{0.7cm}
\setlength{\parindent}{0pt}
\setlength{\parskip}{3pt}
\pagestyle{empty}
\makeatletter
% figure registry (figdefs.tex is generated by the extract step)
\newcommand{\qfigdef}[4]{\@namedef{qf@file@#1}{#2}\@namedef{qf@w@#1}{#3}\@namedef{qf@wide@#1}{#4}}
\newcommand{\qf@check}[1]{\@ifundefined{qf@file@#1}{\PackageError{score2}{Unknown figure id #1}{}}{}}
\newcommand{\qf@img}[4][]{\includegraphics[width=#2\dimexpr\@nameuse{qf@w@#4}pt\relax,max width=#3,max height=0.45\textheight,#1]{\@nameuse{qf@file@#4}}}
% Full-width questions: finalize wraps every question that contains a wide figure in \qfullbegin...\qfullend,
% so the WHOLE question (stem, figures, options) is set across the page in reading order, then the two columns
% resume. The question's head - its number, stem and first figure (\qfullsplit marks the end of the head) -
% is kept on one page, so a number is never left at a page foot with its figure overleaf; the rest may break
% across pages normally, so no big gaps are left. Without a \qfullsplit, the whole question is kept together
% if it is under 45% of a page. Never write these yourself.
\newif\ifqf@full
\newif\ifqf@boxed
\newsavebox{\qf@box}
\newcommand{\qf@open}{\global\qf@fulltrue\par\global\qf@boxedtrue
  \setbox\qf@box\vbox\bgroup\hsize\textwidth\linewidth\textwidth}
\newcommand{\qf@emit}[1]{\par\egroup\global\qf@boxedfalse
  \ifdim\dimexpr\ht\qf@box+\dp\qf@box\relax<#1\textheight\noindent\usebox{\qf@box}\par\else\unvbox\qf@box\fi}
\newcommand{\qfullsplit}{\ifqf@boxed\qf@emit{0.85}\fi}
\newcommand{\qf@close}{\ifqf@boxed\qf@emit{0.45}\fi\global\qf@fullfalse}
\newcommand{\qfullbegin}{\end{multicols}\qf@open}
\newcommand{\qfullend}{\qf@close\begin{multicols}{2}}
\newcommand{\qfullstart}{\qf@open}
\newcommand{\qfulldone}{\qf@close}
% \qfig[scale]{id}: block figure, centred, directly where it is written (size set by the tool).
% Fallback if a wide figure is not inside a full-width question: break out of the columns in place - no floats.
\newcommand{\qfig}[2][1]{\qf@check{#2}%
  \ifqf@full\qf@blk{#1}{#2}\else
  \ifnum\@nameuse{qf@wide@#2}=1 \expandafter\@firstoftwo\else\expandafter\@secondoftwo\fi
  {\end{multicols}\begin{center}\qf@img{#1}{\textwidth}{#2}\end{center}\begin{multicols}{2}}%
  {\qf@blk{#1}{#2}}\fi}
\newcommand{\qf@blk}[2]{\par\smallskip{\centering\qf@img{#1}{\linewidth}{#2}\par}\smallskip}
% \qfigi[scale]{id}: inline figure (inside option lists / tables / minipages); never breaks columns.
\newcommand{\qfigi}[2][1]{\qf@check{#2}\qf@img[valign=c]{#1}{\linewidth}{#2}}
\makeatother
% structure macros - lint reads these to check completeness and figure placement
\newcommand{\papersec}[1]{\par\bigskip{\centering\Large\bfseries PAPER-#1\par}\medskip}
\newcommand{\partsec}[2]{\par\medskip\noindent{\large\bfseries PART-#1: #2}\par\smallskip}
\newcommand{\sect}[1]{\par\smallskip\noindent\textbf{#1}\par}
\newcommand{\qn}[1]{\par\smallskip\noindent\textbf{#1.}~}
\newcommand{\paraq}[2]{\par\smallskip\noindent\textbf{Paragraph for Questions #1 and #2}\par}
\newenvironment{qoptions}{\begin{itemize}[leftmargin=1.8em,itemsep=1pt,topsep=2pt,label={}]}{\end{itemize}}
\newcommand{\opt}[2]{\item[(#1)] #2}
% figure registry (generated by score2_tools.py finalize)
\qfigdef{p1_phy_s1_q01_1}{Nurture (General) - Test 04 (11 Feb 2024)-img/p1_phy_s1_q01_1.png}{122.4}{0}
\qfigdef{p1_phy_s2_q06_1}{Nurture (General) - Test 04 (11 Feb 2024)-img/p1_phy_s2_q06_1.png}{286.9}{1}
\qfigdef{p1_phy_s2_q07_1}{Nurture (General) - Test 04 (11 Feb 2024)-img/p1_phy_s2_q07_1.png}{196.4}{0}
\qfigdef{p1_phy_s2_q08_1}{Nurture (General) - Test 04 (11 Feb 2024)-img/p1_phy_s2_q08_1.png}{134.8}{0}
\qfigdef{p1_phy_s2_q10_1}{Nurture (General) - Test 04 (11 Feb 2024)-img/p1_phy_s2_q10_1.png}{189.4}{0}
\qfigdef{p1_phy_s2_q12_1}{Nurture (General) - Test 04 (11 Feb 2024)-img/p1_phy_s2_q12_1.png}{214.9}{0}
\qfigdef{p1_phy_s3_q16_1}{Nurture (General) - Test 04 (11 Feb 2024)-img/p1_phy_s3_q16_1.png}{154.3}{0}
\qfigdef{p1_chem_s1_q04_A}{Nurture (General) - Test 04 (11 Feb 2024)-img/p1_chem_s1_q04_A.png}{73.1}{0}
\qfigdef{p1_chem_s1_q04_B}{Nurture (General) - Test 04 (11 Feb 2024)-img/p1_chem_s1_q04_B.png}{73.1}{0}
\qfigdef{p1_chem_s1_q04_C}{Nurture (General) - Test 04 (11 Feb 2024)-img/p1_chem_s1_q04_C.png}{73.1}{0}
\qfigdef{p1_chem_s1_q04_D}{Nurture (General) - Test 04 (11 Feb 2024)-img/p1_chem_s1_q04_D.png}{73.1}{0}
\qfigdef{p1_chem_s2_q05_B}{Nurture (General) - Test 04 (11 Feb 2024)-img/p1_chem_s2_q05_B.png}{126.4}{0}
\qfigdef{p1_chem_s2_q05_C}{Nurture (General) - Test 04 (11 Feb 2024)-img/p1_chem_s2_q05_C.png}{90.5}{0}
\qfigdef{p1_chem_s2_q07_A}{Nurture (General) - Test 04 (11 Feb 2024)-img/p1_chem_s2_q07_A.png}{79.6}{0}
\qfigdef{p1_chem_s2_q07_B}{Nurture (General) - Test 04 (11 Feb 2024)-img/p1_chem_s2_q07_B.png}{78.7}{0}
\qfigdef{p1_chem_s2_q07_C}{Nurture (General) - Test 04 (11 Feb 2024)-img/p1_chem_s2_q07_C.png}{79.6}{0}
\qfigdef{p1_chem_s2_q07_D}{Nurture (General) - Test 04 (11 Feb 2024)-img/p1_chem_s2_q07_D.png}{78.7}{0}
\qfigdef{p1_chem_s2_q08_A}{Nurture (General) - Test 04 (11 Feb 2024)-img/p1_chem_s2_q08_A.png}{36.5}{0}
\qfigdef{p1_chem_s2_q08_B}{Nurture (General) - Test 04 (11 Feb 2024)-img/p1_chem_s2_q08_B.png}{44.3}{0}
\qfigdef{p1_chem_s2_q08_C}{Nurture (General) - Test 04 (11 Feb 2024)-img/p1_chem_s2_q08_C.png}{39.5}{0}
\qfigdef{p1_chem_s3_q15_1}{Nurture (General) - Test 04 (11 Feb 2024)-img/p1_chem_s3_q15_1.png}{124.2}{0}
\raggedcolumns
\makeatletter\@beginparpenalty=10000\makeatother
\newcolumntype{L}[1]{>{\raggedright\arraybackslash}p{#1}}
\begin{document}
\begin{center}{\Large\bfseries Nurture (General) --- Test 04 (11 Feb 2024)}\end{center}
\begin{multicols}{2}
\raggedcolumns
\papersec{1}
\partsec{1}{PHYSICS}
\sect{SECTION-I(i)}
\qn{1} A bead c can move freely on a horizontal rod. The bead is connected by blocks B and D by a string as shown in the diagram. If the velocity of block B is $v$. Find the velocity of block D.\par
\qfig{p1_phy_s1_q01_1}
\begin{qoptions}
\opt{A}{$V_D=\dfrac43V$}
\opt{B}{$V_D=\dfrac34V$}
\opt{C}{$V_D=\dfrac35V$}
\opt{D}{$V_D=\dfrac45V$}
\end{qoptions}
\qn{2} The moment of inertia of a semi circular disc of mass M about an axis passing through its centre of mass and perpendicular to the plane is\underline{\hspace{1cm}}
\begin{qoptions}
\opt{A}{$I_{CM}=MR^2\left[\dfrac12-\dfrac4{3\pi^2}\right]$}
\opt{B}{$I_{CM}=MR^2\left[\dfrac12-\dfrac{16}{9\pi^2}\right]$}
\opt{C}{$I_{CM}=MR^2\left[\dfrac12-\dfrac4{\pi^2}\right]$}
\opt{D}{$I_{CM}=\dfrac{MR^2}2+\dfrac{2MR^2}{\pi^2}$}
\end{qoptions}
\qn{3} Two particles of masses m and M are initially at rest at an infinite distance apart. They move towards each other and gain speeds due to gravitational attraction. Find their speeds when the separation between the masses becomes equal to d are
\begin{qoptions}
\opt{A}{$V_1=\sqrt{\dfrac{GM^2}{2d(m+M)}};V_2=\sqrt{\dfrac{Gm^2}{2d(m+M)}}$}
\opt{B}{$V_1=\sqrt{\dfrac{Gm^2}{d(m+M)}};V_2=\sqrt{\dfrac{GM^2}{d(m+M)}}$}
\opt{C}{$V_1=\sqrt{\dfrac{2Gm^2}{d(m+M)}};V_2=\sqrt{\dfrac{2GM^2}{d(m+M)}}$}
\opt{D}{$V_1=\sqrt{\dfrac{2GM^2}{d(m+M)}};V_2=\sqrt{\dfrac{2Gm^2}{d(m+M)}}$}
\end{qoptions}
\qn{4} A chain of mass m and length $\ell$ lies on a rough table. The chain just starts to slip when the overhanging part equals fraction $=\mathrm n\,(0<\mathrm n<1)$ of the chain length. If the chain is slightly disturbed (when overhanging part was fraction $=$ n of chain length) so that it completely slips off the table, what is the work performed by the friction forces
\begin{qoptions}
\opt{A}{$W=-\dfrac12(1-\mathrm n)\mathrm{nmg}\ell$}
\opt{B}{$W=-\dfrac12\mathrm{nmg}\ell$}
\opt{C}{$W=-\dfrac12(1-\mathrm n)\mathrm{nmg}\ell$}
\opt{D}{$W=-\dfrac14(1-\mathrm n)\mathrm{mg}\ell$}
\end{qoptions}
\sect{SECTION-I(ii)}
\qn{5} A car starts moving rectilinearly, first with an acceleration of 5 ms$^{-2}$ (initial velocity zero), then moves uniformly and finally decelerating at the same rate till it stops. The total time of journey is 25 s. The average speed during the interval is $72\,\mathrm{km\,h}^{-1}$. Which of the following statements are correct.
\begin{qoptions}
\opt{A}{Total distance travelled by the particle with uniform acceleration is 62.5 m}
\opt{B}{Total distance travelled by the particle with uniform velocity is 250 m}
\opt{C}{Time duration for which car accelerates is 5 sec}
\opt{D}{The duration for which the car moves with uniform velocity is 15 s}
\end{qoptions}
\qfullbegin
\qn{6} One pan of balance holds a glass of water and stand from which a weight is suspended as in figure, what will not happen to balance's equilibrium if we lengthen the string so as to immerse the weight in the water?\par
\qfig{p1_phy_s2_q06_1}
\qfullsplit
\begin{qoptions}
\opt{A}{Equilibrium will be disrupted since buoyant force acting on weight decreases tension in the string so the pressure of the stand on balance decrease}
\opt{B}{Equilibrium will be disrupted because the weight has been immersed in the water, the water level will increase so that the pressure on the bottom of the glass will increase as well and the left pan will drop down.}
\opt{C}{The equilibrium of the balance will be preserved irrespective of position of weight}
\opt{D}{Instead of lengthening string If you put your finger in the glass of water without touching the walls and the bottom of the glass then also equilibrium will preserved}
\end{qoptions}
\qfullend
\qn{7} In the figure, force $F=\alpha t$ is applied on the block of mass $m_2$. Here $\alpha$ is a constant and t is the time\par
\qfig{p1_phy_s2_q07_1}
\begin{qoptions}
\opt{A}{The acceleration of the blocks are same until $t\le\mu\dfrac{m_1}{\alpha m_2}(m_1+m_2)g$\,.}
\opt{B}{The slope of acceleration-time graph of both the blocks remains same until $t\le\mu\dfrac{m_2}{\alpha m_1}(m_1+m_2)g$}
\opt{C}{Slope of acceleration-time graph of $m_1$ is constant until it separated from $m_2$ whereas for $m_2$ slope changes with time after $t\ge\mu\dfrac{m_2}{\alpha m_1}(m_1+m_2)g$}
\opt{D}{Angle between slope of the graphs after certain interval when relative motion starts is $\tan^{-1}(\alpha/m_2)$}
\end{qoptions}
\qn{8} Consider a uniform disk of mass M and radius R, which is allowed to rotate about a smooth horizontal axis passing through a point on the rim as shown in figure. Initially, the center of mass of disk is in same horizontal level as that of axis of rotation, from this position the disk is released from rest at t = 0. When the disc rotates by angle $\theta$ as shown in figure,\par
\qfig{p1_phy_s2_q08_1}
\begin{qoptions}
\opt{A}{The angular acceleration $\alpha=\left(\dfrac{2g\cos\theta}{3R}\right)$}
\opt{B}{Angular velocity $\omega=\sqrt{\dfrac{4g\sin\theta}{3R}}$}
\opt{C}{Contact force between axle of axis of rotation and disk as a function of angular displacement is given by $N=\dfrac{Mg}3\sqrt{1+48\sin^2\theta}$}
\opt{D}{Acceleration of center of mass $a_{cm}=\dfrac{2g}3\sqrt{1+3\sin^2\theta}$}
\end{qoptions}
\qn{9} A man is standing on a road and observes that rain is falling at angle 45$^\circ$ with the vertical. The man starts running on the road with constant acceleration 0.5 m/s$^2$. After a certain time from the start of the motion, it appears to him that rain is still falling at angle 45$^\circ$ with the vertical, with speed $2\sqrt2$ m / s. Motion of the man is in the same vertical plane in which the rain is falling. Then which of the following statement(s) are true
\begin{qoptions}
\opt{A}{It is not possible}
\opt{B}{Speed of the rain relative to the ground is 2 m/s}
\opt{C}{Speed of the man when he finds rain to be falling at angle 45$^\circ$ with the vertical, is 4 m/s}
\opt{D}{The man has travelled a distance 16 m on the road by the time he again finds rain to be falling at angle 45$^\circ$}
\end{qoptions}
\qn{10} Two spheres are moving towards each other. Both have same radius but their masses are 2 kg and 4 kg. If the velocities are 4 m/s and 2 m/s respectively and coefficient of restitution is e = 1/3, then\par
\qfig{p1_phy_s2_q10_1}
\begin{qoptions}
\opt{A}{Maximum potential energy of deformation 18 joule}
\opt{B}{Final velocities along line of impact $v_1=\dfrac{-2}{\sqrt3}$m / s and $v_2=\dfrac1{\sqrt3}$m / s}
\opt{C}{Impulse of deformation $=$ Impulse of reformation}
\opt{D}{Law of conservation of linear momentum along line of motion is applicable and law of conservation energy is also applicable}
\end{qoptions}
\qn{11} A block with mass M attached to a horizontal spring with force constant k is moving with simple harmonic motion having amplitude $A_1$. At the instant when the block passes through its equilibrium position a lump of putty with mass m is dropped vertically on the block from a very small height and sticks to it.
\begin{qoptions}
\opt{A}{The new amplitude is $A_2=A_1\sqrt{\dfrac M{M+m}}$}
\opt{B}{Time period will be $T_2=2\pi\sqrt{\dfrac{M+m}k}$}
\opt{C}{For the case in which the putty is dropped on the block when it is at one end of its path, the new amplitude $A_2=A_1\sqrt{\dfrac m{m+M}}$}
\opt{D}{For the case in which the putty is dropped on the block when it is at one end of its path, the new Time period will be $T_2=2\pi\sqrt{\dfrac{m+M}{2K}}$}
\end{qoptions}
\qn{12} In the given figure there are two blocks A and B. They are given initial velocities as shown. The spring is in natural length initially.\par
\qfig{p1_phy_s2_q12_1}
\begin{qoptions}
\opt{A}{The maximum stretch in the spring is $V_0\sqrt{\dfrac{6m}k}$}
\opt{B}{Velocity of center of mass is $V_0$}
\opt{C}{In center of mass frame velocity of A is $2V_0$}
\opt{D}{In zero momentum frame of reference velocity of B is $V_0$}
\end{qoptions}
\sect{SECTION-II}
\qn{13} A slowly moving flat car is 12.0 m long passing a point at straight road at 10 kmph. A boy beside the road near that point tosses rocks on to the moving flat car at the rate of one per second. How many rocks will fall onto that if the train begins to accelerate at $0.5\,\mathrm{m/s^2}$ just as the first rock hits the car?
\qn{14} A closed organ pipe and an open organ pipe of same length produce 2 beats when they are set into vibrations simultaneously in their fundamental mode. The length of open organ pipe is now halved and of closed organ pipe is doubled, the number of beats produced when they vibrate simultaneously in their fundamental mode will be \underline{\hspace{1cm}}
\qn{15} A particle of mass 2kg is moving under the influence of a force which always acts towards the center and whose potential energy is given by $U(r)=2r^3j$. If the body is moving in a circular orbit of radius 5m, then find its energy in J.
\qn{16} For the figure shown, the wedge $+$ particle is resting at a vertical wall. The particle of mass $m=\dfrac M2$ is released from rest from the position shown. The wedge has mass M. After crossing point A, the maximum height above A reached by the particle on the circular groove is 2R/n, find n. Point A is lowest print on circular groove.\par
\qfig{p1_phy_s3_q16_1}
\qn{17} A metallic sphere of radius $1.0\times10^{-3}$m and density $1.0\times10^4$kg / m$^3$ enters a tank of water, after a free fall through a distance of h in the earth's gravitational field. If its velocity remains unchanged after entering water, determine the value of h(in m). Given: coefficient of viscosity of water $=1.0\times10^{-3}$N $-$ s / m$^2$, $g=10$ m / s$^2$ and density of water $=1.0\times10^3$kg / m$^3$.
\qn{18} 1 mole of an ideal monoatomic gas is expanded till the temperature of the gas is doubled under the process $V^2T$ = constant. The initial temperature of the gas is 400 K. In terms of R, the total work done in the process is -nR find n
\partsec{2}{CHEMISTRY}
\sect{SECTION-I(i)}
\qn{1} When Cu$_2$S is converted into Cu$^{2+}$ \& SO$_2$, then equivalent weight of Cu$_2$S will be (M = mol.wt. of Cu$_2$S)
\begin{qoptions}
\opt{A}{M}
\opt{B}{$\dfrac M2$}
\opt{C}{$\dfrac M4$}
\opt{D}{$\dfrac M8$}
\end{qoptions}
\qn{2} For the equilibrium H$_2$O(s)$\rightleftharpoons$H$_2$O(l)\par
Which of the following statements is true ?
\begin{qoptions}
\opt{A}{The pressure changes do not effect the equilibrium}
\opt{B}{More of ice melts, if pressure on the system is increased}
\opt{C}{More of liquid freezes, if pressure on the system is increased}
\opt{D}{Less of ice melts, if pressure on the system is increased}
\end{qoptions}
\qn{3} Molecular shapes of ClF$_3$, I$_3^-$ and XeO$_3$ respectively are
\begin{qoptions}
\opt{A}{T-shape, Linear, Pyramidal}
\opt{B}{Planar, Linear, Tetrahedral}
\opt{C}{T-shape, Planar, Pyramidal}
\opt{D}{Trigonal bipyramidal, Linear, Tetrahedral}
\end{qoptions}
\qn{4} Graph of incident frequency with stopping potential in photoelectric effect is
\begin{qoptions}
\opt{A}{\qfigi{p1_chem_s1_q04_A}}
\opt{B}{\qfigi{p1_chem_s1_q04_B}}
\opt{C}{\qfigi{p1_chem_s1_q04_C}}
\opt{D}{\qfigi{p1_chem_s1_q04_D}}
\end{qoptions}
\sect{SECTION-I(ii)}
\qn{5} Which of the following correctly represents the stability of reactive intermediate ?
\begin{qoptions}
\opt{A}{$\mathrm{CH_3-\overset{+}{C}H_2<CH_3O\overset{\oplus}{C}H_2}$}
\opt{B}{\qfigi{p1_chem_s2_q05_B}}
\opt{C}{\qfigi{p1_chem_s2_q05_C}}
\opt{D}{$\mathrm{(CH_3)_3\overset{\oplus}{C}<CH_3\overset{\oplus}{C}H_2}$}
\end{qoptions}
\qn{6} In which of the following cases product will contain more number of carbon atoms than do present in reactant molecule ?
\begin{qoptions}
\opt{A}{RCH$=$CHR$'$ $+$ H$_2\longrightarrow$ A}
\opt{B}{RCOOK $+$ OH$^-\xrightarrow[\text{CaO}]{\text{heat}}$ B}
\opt{C}{RCOOK $+$ H$_2$O $\xrightarrow{\text{Electrolysis}}$ C}
\opt{D}{RX $+$ Na $\xrightarrow{\text{dry ether}}$ D}
\end{qoptions}
\qn{7} Among the following, reaction(s) which gives (give) tert-butyl benzene as the major product is(are)
\begin{qoptions}
\opt{A}{\qfigi{p1_chem_s2_q07_A}}
\opt{B}{\qfigi{p1_chem_s2_q07_B}}
\opt{C}{\qfigi{p1_chem_s2_q07_C}}
\opt{D}{\qfigi{p1_chem_s2_q07_D}}
\end{qoptions}
\qn{8} Which of the following will have C$=$O bond almost similar to C$-$O bond length ?
\begin{qoptions}
\opt{A}{\qfigi{p1_chem_s2_q08_A}}
\opt{B}{\qfigi{p1_chem_s2_q08_B}}
\opt{C}{\qfigi{p1_chem_s2_q08_C}}
\opt{D}{O$=$C$=$O}
\end{qoptions}
\qn{9} Among the following, the state function(s) is (are)
\begin{qoptions}
\opt{A}{Internal energy}
\opt{B}{Irreversible expansion work}
\opt{C}{Reversible expansion work}
\opt{D}{Molar enthalpy}
\end{qoptions}
\qn{10} Aqueous solutions of HNO$_3$, KOH, CH$_3$COOH, and CH$_3$COONa of identical concentrations are provided. The pair(s) of solutions which form a buffer upon mixing is (are)
\begin{qoptions}
\opt{A}{HNO$_3$ and CH$_3$COOH}
\opt{B}{KOH and CH$_3$COONa}
\opt{C}{HNO$_3$ and CH$_3$COONa}
\opt{D}{CH$_3$COOH and CH$_3$COONa}
\end{qoptions}
\qn{11} Choose the process which is/are endothermic.
\begin{qoptions}
\opt{A}{O$^{-1}$ + e$^-\longrightarrow$ O$^{-2}$}
\opt{B}{Ar$\longrightarrow$ Ar$^+_{\,}$ + e$^-$}
\opt{C}{Ar + e$^-\longrightarrow$ Ar$^-$}
\opt{D}{H + e$^-\longrightarrow$ H$^-$}
\end{qoptions}
\qn{12} Among the following, the correct statement(s) is (are)
\begin{qoptions}
\opt{A}{AlCl$_3$ has the three-centre two-electron bonds in its dimeric structure}
\opt{B}{BH$_3$ has the three-centre two-electron bonds in its dimeric structure}
\opt{C}{Al(CH$_3$)$_3$ has the three-centre two-electron bonds in its dimeric structure}
\opt{D}{The Lewis acidity of BCl$_3$ is greater than that of AlCl$_3$}
\end{qoptions}
\sect{SECTION-II}
\qn{13} Amongst the following, the total number of compounds whose aqueous solution turns red litmus paper blue is\par
KCN \quad K$_2$SO$_4$ \quad (NH$_4$)$_2$C$_2$O$_4$ \quad NaCl \quad Zn(NO$_3$)$_2$ \quad FeCl$_3$ \quad K$_2$CO$_3$ \quad NH$_4$NO$_3$ \quad LiCN
\qn{14} The value of n in the molecular formula Be$_n$Al$_2$Si$_6$O$_{18}$ is
\qn{15} The number of resonance structures for N is\par
\qfig{p1_chem_s3_q15_1}
\qn{16} The maximum number of electrons that can have principal quantum number, n = 3, and spin quantum number, m$_s=-\dfrac12$, is
\qn{17} In the reaction VO $+$ Fe$_2$O$_3\to$ FeO $+$ V$_2$O$_5$, what is the n-factor for V$_2$O$_5$ ?
\qn{18} Among H$_2$, He$_2^+$, Li$_2$, Be$_2$, B$_2$, C$_2$, N$_2$, O$_2^-$, and F$_2$, the number of diamagnetic species is (At. numbers : H = 1, He = 2, Li = 3, Be = 4, B = 5, C = 6, N = 7, O = 8, F = 9)
\partsec{3}{MATHEMATICS}
\sect{SECTION-I(i)}
\qn{1} The sum $\dfrac{19}{1.2.3}\dfrac14+\dfrac{28}{2.3.4}\dfrac18+\dfrac{39}{3.4.5}\dfrac1{16}+\dfrac{52}{4.5.6}\dfrac1{32}+\ldots+$ upto infinite terms is equal to :
\begin{qoptions}
\opt{A}{$\dfrac12$}
\opt{B}{1}
\opt{C}{2}
\opt{D}{4}
\end{qoptions}
\qn{2} If $|z-2-i|=|z|\sin\!\left(\dfrac\pi4-\arg z\right)$, then locus of z is/an:
\begin{qoptions}
\opt{A}{Ellipse}
\opt{B}{Circle}
\opt{C}{Parabola}
\opt{D}{Pair of straight lines}
\end{qoptions}
\qn{3} Number of positive integral solutions satisfying the equation $(x_1+x_2+x_3)(y_1+y_2)=77$, is:
\begin{qoptions}
\opt{A}{150}
\opt{B}{270}
\opt{C}{420}
\opt{D}{1024}
\end{qoptions}
\qn{4} a, b, c, d are distinct integers such that $(x-a)(x-b)(x-c)(x-4)=4$ has an integral root r. Then $a+b+c+d$ is equal to :
\begin{qoptions}
\opt{A}{r}
\opt{B}{2r}
\opt{C}{3r}
\opt{D}{4r}
\end{qoptions}
\sect{SECTION-I(ii)}
\qn{5} The sequence $\{a_n\}$, $n\in N$ satisfies $a_1=1$ and $5^{a_{n+1}-a_n}=1+\dfrac1{n+\frac23}$. Then :\par
(where [.] denotes greatest integer function)
\begin{qoptions}
\opt{A}{$[a_{501}]=3$}
\opt{B}{$[a_{207}]=3$}
\opt{C}{$[a_{223}]=4$}
\opt{D}{$[a_{625}]=4$}
\end{qoptions}
\qn{6} Which of the following is/are positive ?
\begin{qoptions}
\opt{A}{$\log_{\sin1}\tan1$}
\opt{B}{$\log_{\cos1}(1+\tan3)$}
\opt{C}{$\log_{\log_{10}3^{\,}}(\cos\theta+\sec\theta)$}
\opt{D}{$\log_{\tan15^\circ}(2\sin18^\circ)$}
\end{qoptions}
\qn{7} Let $S_1=\displaystyle\sum_{r=0}^{n}\left({}^{2n+1}C_{2r}\right)^2$ and $S_2=\displaystyle\sum_{r=0}^{n}\left({}^{2n+1}C_{2r+1}\right)^2$, then
\begin{qoptions}
\opt{A}{$S_1=\dfrac12\left({}^{4n+2}C_{2n}+(-1)^n\,{}^{2n+1}C_n\right)$}
\opt{B}{$S_2=\dfrac12\left({}^{4n+2}C_{2n}-(-1)^n\,{}^{2n+1}C_n\right)$}
\opt{C}{$S_1=\dfrac12\left({}^{4n+2}C_{2n+1}\right)$}
\opt{D}{$S_2=\dfrac12\left({}^{4n+1}C_{2n+1}\right)$}
\end{qoptions}
\qn{8} If equation of tangent at P, Q and vertex A of a parabola are $3x+4y-7=0$, $2x+3y-10=0$ and $x-y=0$ respectively, then :
\begin{qoptions}
\opt{A}{Focus is (4, 5)}
\opt{B}{Length of latus rectum is $2\sqrt2$}
\opt{C}{Axis is $x+y-9=0$}
\opt{D}{Vertex is $\left(\dfrac92,\dfrac92\right)$}
\end{qoptions}
\qn{9} Equation of line through a and ib such that $a,b\in R$ and $a,b\ne0$ is
\begin{qoptions}
\opt{A}{$z\left(\dfrac1{2a}+\dfrac1{2ib}\right)+\overline z\left(\dfrac1{2a}-\dfrac1{2ib}\right)=1$}
\opt{B}{$z\left(\dfrac1{2a}+\dfrac i{2b}\right)+\overline z\left(\dfrac1{2a}+\dfrac i{2b}\right)=1$}
\opt{C}{$z\left(\dfrac1{2a}-\dfrac i{2b}\right)+\overline z\left(\dfrac1{2a}-\dfrac i{2b}\right)=1$}
\opt{D}{$z\left(\dfrac1{2a}-\dfrac i{2b}\right)+\overline z\left(\dfrac1{2a}+\dfrac i{2b}\right)=1$}
\end{qoptions}
\qn{10} Number of ways in which the letters of the word `B U L B U L' can be arranged in a line in any order is also equal to the
\begin{qoptions}
\opt{A}{Number of ways in which 2 alike Apples and 4 alike Mangoes can be distributed in 3 children so that each child receives any number of fruits.}
\opt{B}{Number of ways in which 6 different books can be tied up into 3 bundles, if each bundle is to have equal number of books}
\opt{C}{Coefficient of $x^2y^2z^2$ in the expansion of $(x+y+z)^6$}
\opt{D}{Number of ways in which 6 different prizes can be distributed equally in three children.}
\end{qoptions}
\qn{11} Consider the equation $\sqrt{x+\sqrt{2x-1}}+\sqrt{x-\sqrt{2x-1}}=A$
\begin{qoptions}
\opt{A}{For $A=\sqrt2$, $x\in\left[\dfrac12,1\right]$}
\opt{B}{For $A=\sqrt2$, $x\in\left[0,\dfrac12\right]$}
\opt{C}{For $A=1$, $x\in\phi$}
\opt{D}{For $A=2$, $x=\dfrac32$}
\end{qoptions}
\qn{12} The co-ordinates (2, 3) and (1, 5) are the foci of an ellipse which passes through the origin, then the equation of
\begin{qoptions}
\opt{A}{Tangent at the origin is $(3\sqrt2-5)x+(1-2\sqrt2)y=0$}
\opt{B}{Tangent at the origin is $(3\sqrt2+5)x-(1+2\sqrt2)y=0$}
\opt{C}{Normal at the origin is $(3\sqrt2+5)x-(2\sqrt2+1)y=0$}
\opt{D}{Normal at the origin is $(3\sqrt2-5)x+(1-2\sqrt2)y=0$}
\end{qoptions}
\sect{SECTION-II}
\qn{13} Transverse and conjugate axes of a rectangular hyperbola are along X-axis and Y-axis respectively and the distance between the foci is $10\sqrt{14}$. Number of the points (x, y) on the curve such that x and y are positive integers, is equal to \underline{\hspace{1cm}}.
\qn{14} $96\sqrt3\,\sin\dfrac\pi{48}\cos\dfrac\pi{48}\cos\dfrac\pi{24}\cos\dfrac\pi{12}\cos\dfrac\pi6$ has the value $=$ \ldots\ldots\ldots\ldots
\qn{15} The locus of the centre of a circle which touches the circles $|z-z_1|=2$ and $|z-z_2|=4$, externally is a hyperbola if $|z_1-z_2|$, is greater than \underline{\hspace{1cm}}
\qn{16} If the circles $x^2+y^2+(3+\sin\beta)x+(2\cos\alpha)y=0$ and $x^2+y^2+(2\cos\alpha)x+2cy=0$ touch each other, then the maximum value of c is \underline{\hspace{1cm}}.
\qn{17} Evaluate the sum $\sqrt{1-\dfrac12\cdot\sqrt{1\cdot3}}+\sqrt{2-\dfrac12\cdot\sqrt{3\cdot5}}$ $+\sqrt{3-\dfrac12\cdot\sqrt{5\cdot7}}$\par $+\cdots+\sqrt{40-\dfrac12\cdot\sqrt{79\cdot81}}$
\qn{18} Three positive reals x, y, z satisfy\par
$x^2+y^2=3^2$\par
$y^2+yz+z^2=4^2$\par
$x^2+\sqrt3xz+z^2=5^2$\par
Find the value of $2xy+xz+\sqrt3yz$
\end{multicols}
\end{document}
